\chapter{使用过程中的常见问题}

\section{为什么模板能打开但不能编译}
最常见的原因并不是模板本身损坏，而是某些必要文件缺失，例如章节文件被删除、图片路径写错、参考文献数据库不存在，或者正文中出现了不完整的环境。若遇到这类问题，建议优先查看 `main.log` 最后几十行的报错信息，而不是直接从 PDF 结果猜原因。

\section{为什么版式和样张不完全一致}
LaTeX 在图表浮动、分页和长段落换行方面，会和 Word 排版存在一些天然差异。因此模板能做到的是在页边距、字号、标题、目录、图题、表题和常见正文样式上尽量贴近手册样张，而不是保证每一页都与样张逐像素一致。最终提交前，仍建议逐页检查版面。

\section{哪些内容适合在类文件中改}
若要修改页边距、标题格式、目录样式、正文行距、图表题注或页码格式，建议统一在 `gdthesis.cls` 中调整。这样能够保证整篇论文保持一致，也便于后续维护。如果只是修改论文内容，则不建议频繁改动类文件。

\section{给使用者的最后建议}
最稳妥的写作方式通常是先把结构写完整，再逐步替换每一章内容，而不是从零开始搭整个模板。也就是说，先保证“能编译、能出目录、能出章节、能出参考文献”，然后再去追求更细的版式微调。这样效率会更高，也更不容易在后期因为结构性问题返工。
